
\documentclass[11pt]{article}
\usepackage{series}
\hypersetup{
  pdftitle={Finite-Time Scattering Windows and Asymptotic Delta Limits},
  pdfauthor={Peter Nero},
  pdfcreator={LaTeX}
}

\title{Finite-Time Scattering Windows and Asymptotic Delta Limits\\
\large Standard Fourier Mathematics and the MTT Source Boundary}
\author{Peter Nero}
\date{July 2026, Version 1}

\begin{document}
\maketitle

\begin{abstract}
The exact energy--momentum conserving delta functions of the scattering matrix,
\[
(2\pi)^4\delta^{(4)}(p_f-p_i),
\]
are usually introduced as consequences of translation invariance and asymptotic scattering theory.  This paper continues the sequence interpreting Dirac deltas as singular shadows of admissible projection.  The previous bookkeeping paper showed that momentum-conservation deltas arise as infinite-support limits of finite interaction-window Fourier kernels.  Here we focus on scattering itself: finite-time transition amplitudes do not produce exact energy deltas, but finite-width sinc or window kernels.  The exact S-matrix delta appears only after the asymptotic idealization of infinite observation time, infinite spatial support, and stable in/out sectors.

We prove a standard finite-time transition theorem: for a constant perturbation acting during a time window of length \(T\), the first-order transition amplitude contains
\[
W_T(\Delta E)=\int_{-T/2}^{T/2}e^{i\Delta E t}\,dt
=2\frac{\sin(\Delta E T/2)}{\Delta E},
\]
and the normalized squared kernel converges distributionally to \(2\pi\delta(\Delta E)\).  More generally, smooth time windows yield approximate energy-conservation kernels through their Fourier transforms, with normalization fixed by Plancherel.  The transition-rate limit reproduces the usual Fermi-golden-rule factor, and wave-packet scattering shows how plane-wave deltas are always integrated against finite profiles in physical amplitudes.  We then combine finite time and finite spatial support to obtain finite energy--momentum bookkeeping kernels replacing the exact S-matrix conservation delta.

The MTT interpretation is conditional.  These standard window kernels may serve as
bookkeeping outputs if selected MTT dynamics derives the preparation, switching, detector,
or finite-time profile.  This paper does not derive such a profile, and finite experimental
duration must not be identified with a fundamental admissibility width.  Exact conservation
continues to follow from the symmetry assumptions of the closed theory.
\end{abstract}

\section*{Version 1 Revision Note}
\begin{description}
\item[Supersedes] The unversioned April 2026 manuscript.
\item[Reason] The finite-window limits were standard and correct, but the earlier MTT
language could make an externally chosen observation window sound geometrically selected.
\item[Resolution] This version separates exact symmetry conservation, finite experimental
resolution, and a possible fundamental MTT width.  Only the first two are established here.
\item[Retained result] The sinc, smooth-window, Plancherel, golden-rule, and spacetime
distributional limits are unchanged.
\item[Remaining boundary] Selected dynamics must emit any claimed MTT transition window.
\end{description}

\tableofcontents

\section{Purpose and claim discipline}

This paper is the scattering sequel to the delta/projection series.  Its purpose is narrow.

We do not claim to reconstruct full scattering theory from first principles in this paper.  We also do not claim that energy or momentum conservation is merely approximate in a closed system.  Rather, we prove and interpret a standard fact:

\begin{quote}
Exact S-matrix delta functions arise only after idealizing finite transition windows into infinite asymptotic support.
\end{quote}

The mathematical statement is elementary Fourier analysis.  The MTT interpretation is an
additional source question: a kernel becomes an MTT output only if selected dynamics emits
the corresponding preparation, interaction, or detector window.

\begin{center}
\fbox{\begin{minipage}{0.86\linewidth}
\textbf{Core claim.}  The S-matrix delta
\[
(2\pi)^4\delta^{(4)}(p_f-p_i)
\]
is the asymptotic, infinite-support limit of finite energy--momentum bookkeeping kernels generated by finite transition windows.
\end{minipage}}
\end{center}

\section{Three distinct notions of width}

The same formula can occur in three logically different settings:
\begin{enumerate}
\item a switching or observation time chosen by an experiment;
\item finite-volume or finite-duration support that breaks exact translation invariance in an
effective description;
\item a fundamental admissibility width selected by an underlying theory.
\end{enumerate}
The first two are ordinary scattering inputs.  The third is not established by the Fourier
limit.  In particular, a measured duration \(T\) does not determine an MTT coherence scale.

\section{Fourier conventions}

We use the following convention on time:
\[
\widehat w(\omega)=\int_{\mathbb R} w(t)e^{i\omega t}\,dt.
\]
With this convention,
\[
\lim_{T\to\infty}\int_{-T/2}^{T/2}e^{i\omega t}\,dt
=2\pi\delta(\omega)
\]
in the sense of distributions.

In spacetime dimension \(d+1\), for a window \(w(x)\) with \(x=(t,\mathbf x)\), we write
\[
\widehat w(q)=\int w(x)e^{iq\cdot x}\,d^{d+1}x,
\]
with the corresponding limiting convention
\[
\widehat w_R(q)\longrightarrow (2\pi)^{d+1}\delta^{(d+1)}(q)
\]
for approximate-identity windows whose support scale tends to infinity.

\section{Finite-time transition amplitudes}

Let \(H=H_0+V\) and suppose \(|i\rangle,|f\rangle\) are eigenstates of \(H_0\) with
\[
H_0|i\rangle=E_i|i\rangle,\qquad H_0|f\rangle=E_f|f\rangle.
\]
In first-order time-dependent perturbation theory, with perturbation \(V\) active only through a time window \(w_T(t)\), the transition amplitude is
\[
\mathcal A_{fi}^{(1)}(T)
=
-i\,V_{fi}\int_{\mathbb R} w_T(t)e^{i(E_f-E_i)t}\,dt.
\]
Writing
\[
\Delta E:=E_f-E_i,
\]
we have
\[
\mathcal A_{fi}^{(1)}(T)
=
-i\,V_{fi}\,\widehat w_T(\Delta E).
\]

Thus the finite-time transition does not contain \(\delta(\Delta E)\).  It contains the finite bookkeeping kernel \(\widehat w_T(\Delta E)\).

\section{Box window and the sinc kernel}

For the sharp observation window
\[
w_T(t)=\mathbf 1_{[-T/2,T/2]}(t),
\]
the kernel is
\[
W_T(\omega)=\widehat w_T(\omega)
=
\int_{-T/2}^{T/2}e^{i\omega t}\,dt
=
2\frac{\sin(\omega T/2)}{\omega}.
\]

This is a sinc-type energy bookkeeping kernel.  It is sharply peaked near \(\omega=0\), but for finite \(T\) it is not a delta.

\begin{theorem}[Finite-time energy kernel]
Let
\[
D_T(\omega):=\frac{1}{2\pi T}|W_T(\omega)|^2
=
\frac{1}{2\pi T}\left(\frac{2\sin(\omega T/2)}{\omega}\right)^2.
\]
Then
\[
D_T(\omega)\to \delta(\omega)
\]
in the sense of distributions as \(T\to\infty\).
\end{theorem}

\begin{proof}
Let \(f\) be a Schwartz test function.  With \(u=\omega T/2\), one obtains
\[
\int_{\mathbb R}D_T(\omega)f(\omega)\,d\omega
=
\int_{\mathbb R}\frac{1}{\pi}\left(\frac{\sin u}{u}\right)^2
f\!\left(\frac{2u}{T}\right)\,du.
\]
Since
\[
\int_{\mathbb R}\frac{1}{\pi}\left(\frac{\sin u}{u}\right)^2du=1,
\]
and \(f(2u/T)\to f(0)\) pointwise, dominated-convergence/standard approximate-identity arguments give
\[
\int_{\mathbb R}D_T(\omega)f(\omega)\,d\omega\to f(0).
\]
Thus \(D_T\to\delta\) distributionally.
\end{proof}

The unsquared amplitude kernel itself satisfies
\[
W_T(\omega)\to 2\pi\delta(\omega)
\]
distributionally.  The squared, normalized kernel \(D_T\) is the probability-rate version.


\begin{center}
\begin{tabular}{lll}
\hline
Object & Finite-time expression & Asymptotic limit \\
\hline
Amplitude kernel & \(W_T(\omega)\) & \(2\pi\delta(\omega)\) \\
Rate kernel & \(\frac{1}{T}|W_T(\omega)|^2\) & \(2\pi\delta(\omega)\) \\
Normalized rate kernel & \(\frac{1}{2\pi T}|W_T(\omega)|^2\) & \(\delta(\omega)\) \\
\hline
\end{tabular}
\end{center}

The distinction is important.  Amplitudes contain a distributional delta in the infinite-time limit.  Probabilities contain squared kernels, and the meaningful object is the rate-normalized approximate identity.

\section{Fermi's golden rule as the rate limit}

The first-order transition probability is
\[
P_{fi}(T)=|V_{fi}|^2 |W_T(\Delta E)|^2.
\]
For transitions into a continuum of final states with density \(\rho(E_f)\), the total probability is
\[
P_i(T)=\int |V_{fi}|^2 |W_T(E_f-E_i)|^2\rho(E_f)\,dE_f.
\]
Using the distributional limit
\[
\frac{1}{T}|W_T(\omega)|^2\to 2\pi\delta(\omega),
\]
one obtains the transition rate
\[
\Gamma_i
=
\lim_{T\to\infty}\frac{P_i(T)}{T}
=
2\pi |V_{fi}|^2\rho(E_i),
\]
in the usual idealized form.

Thus Fermi's golden rule already displays the core pattern of this paper:

\[
\boxed{\text{finite-time transition kernel}\quad\longrightarrow\quad\text{asymptotic energy delta}.}
\]

\section{Smooth time windows}

The box window is useful because it gives the familiar sinc kernel, but the conclusion does not depend on the discontinuity of the box.  Let
\[
w_T(t)=w(t/T),
\]
with \(w\in C_c^\infty(\mathbb R)\) nonzero.  Then
\[
\widehat w_T(\omega)=T\widehat w(T\omega).
\]
With our Fourier convention, Plancherel gives
\[
\|\widehat w\|_{L^2(\omega)}^2=2\pi\|w\|_{L^2(t)}^2.
\]
Hence the correctly normalized smooth-window rate kernel is
\[
D_T^w(\omega)
=
\frac{|\widehat w_T(\omega)|^2}{2\pi T\|w\|_{L^2}^2}
=
\frac{T|\widehat w(T\omega)|^2}{2\pi \|w\|_{L^2}^2}.
\]
It has unit integral and forms an approximate identity:
\[
D_T^w(\omega)\to \delta(\omega)
\]
distributionally.

Indeed, for a Schwartz test function \(f\),
\[
\int_{\mathbb R}D_T^w(\omega)f(\omega)\,d\omega
=
\int_{\mathbb R}
\frac{|\widehat w(u)|^2}{2\pi\|w\|_{L^2}^2}
f(u/T)\,du
\to f(0).
\]
The constants change if one uses a different Fourier convention, but the structural result is invariant: finite time windows give finite-width energy kernels, and exact deltas appear only in the infinite-time limit.

This shows that the exact energy delta is not tied to a discontinuous box window.  It is the universal infinite-time limit of finite admissible time windows.

\section{Wave-packet scattering}

Plane-wave scattering states are themselves idealizations: they are infinitely extended and exactly sharp in momentum.  A finite incoming packet has the form
\[
|\psi_i\rangle=\int dp\, f_i(p)|p\rangle,
\qquad
|\psi_f\rangle=\int dp\, f_f(p)|p\rangle,
\]
with normalized packet profiles \(f_i,f_f\).  When the plane-wave S-matrix contains a conservation factor,
\[
(2\pi)^{d+1}\delta^{(d+1)}(p_f-p_i),
\]
the packet amplitude contains this distribution only under integration against the packet profiles:
\[
\mathcal A_{\psi_f\psi_i}
=
\int dp_f\,dp_i\,
f_f^\ast(p_f)f_i(p_i)
(2\pi)^{d+1}\delta^{(d+1)}(p_f-p_i)\mathcal M(p_f,p_i).
\]
For finite packets and finite interaction windows, the delta is replaced by a finite overlap kernel,
\[
\delta^{(d+1)}(p_f-p_i)
\quad\leadsto\quad
K_{T,R}(p_f-p_i),
\]
where \(K_{T,R}\) is the Fourier transform of the spacetime transition window.  The physically observed amplitude is therefore a packet-weighted overlap, not a naked delta.  The exact delta belongs to the plane-wave, infinite-support idealization.

In MTT terms, wave packets are closer to admissible finite bookkeeping than plane waves.  Plane waves expose the singular limit cleanly, but finite packets better represent the finite-resolution transition structure.

\section{LSZ, asymptotic regimes, and admissible scattering}

The exact S-matrix is not available in arbitrary dynamical regimes.  It requires stable
asymptotic sectors in which in/out states can be defined and compared.  In ordinary QFT this
is the role of Haag--Ruelle/LSZ-type scattering theory in regimes with appropriate mass gaps,
stability, and asymptotic separation \cite{LSZ1955,Ruelle1962}.

The present paper does not derive LSZ or Haag--Ruelle theory.  It uses their standard lesson as a domain condition:

\begin{quote}
S-matrix deltas are meaningful only when the theory admits stable asymptotic in/out sectors.
\end{quote}

In MTT language, this means that an S-matrix description is an admissible scattering encoding, not a universal description of all processes.  In time-dependent backgrounds, finite systems, measurement contexts, or strongly interacting non-asymptotic regimes, local observables or in-in quantities may be the correct objects, and exact conservation deltas should be replaced by finite bookkeeping kernels.


\section{Finite spatial support and full energy--momentum kernels}

Let an interaction be supported by a spacetime window
\[
w_{T,R}(t,\mathbf x)
=
w_T(t)v_R(\mathbf x).
\]
For an interaction with total energy--momentum mismatch
\[
q=(\Delta E,\Delta \mathbf p),
\]
the vertex factor contains
\[
\widehat w_{T,R}(q)
=
\widehat w_T(\Delta E)\widehat v_R(\Delta \mathbf p).
\]
For finite \(T,R\), this is a finite-width energy--momentum bookkeeping kernel.  In the asymptotic limit,
\[
\widehat w_{T,R}(q)
\to
(2\pi)^{d+1}\delta(\Delta E)\delta^{(d)}(\Delta\mathbf p).
\]

Hence the exact S-matrix conservation factor is the infinite-time, infinite-volume limit of finite transition support:
\[
(2\pi)^{d+1}\delta^{(d+1)}(q)
=
\lim_{T,R\to\infty}\widehat w_{T,R}(q).
\]

\section{S-matrix deltas and asymptotic idealization}

In standard scattering theory, one writes
\[
\langle f|S|i\rangle
=
\delta_{fi}
+
i(2\pi)^4\delta^{(4)}(p_f-p_i)\mathcal M_{fi}.
\]
The delta is not produced by a finite laboratory process.  It belongs to the asymptotic idealization in which:

\begin{enumerate}
\item in/out states are stable and well separated;
\item the interaction region is effectively isolated;
\item time support has been idealized to infinity;
\item spatial support has been idealized to exact translation invariance;
\item bookkeeping mismatch has been forced to zero width.
\end{enumerate}

This is not a criticism of the S-matrix.  It is a clarification of its domain.  The S-matrix is the correct object when admissible scattering regimes exist.  In finite-time or time-dependent regimes, local or in-in observables are usually more appropriate, and conservation deltas are replaced by finite kernels.

\section{MTT interpretation}

The MTT reading is that exact scattering deltas belong to the circle/bookkeeping side of the triadic carrier.  They express exact closure of the energy--momentum ledger after the finite transition window has been idealized away.

The finite kernel
\[
\widehat w_{T,R}(q)
\]
is the native object.  It measures how strongly a process with mismatch \(q\) is supported by the finite admissible transition region.  The delta arises only when the bookkeeping window becomes infinitely extended and the mismatch width collapses to zero.

\[
\boxed{
\text{S-matrix delta}
=
\text{singular asymptotic shadow of finite bookkeeping closure}.
}
\]

This completes the circle-sector counterpart to the previous papers:

\begin{center}
\begin{tabular}{lll}
\hline
Paper & Delta role & MTT sector emphasis \\
\hline
Gauge fixing & Representative selection & Lens redundancy \\
Measurement & Survivor-basin selection & Nil stabilization \\
Momentum/scattering & Conservation bookkeeping & Circle closure \\
\hline
\end{tabular}
\end{center}

\section{Conservation is not merely approximate}

A possible misunderstanding should be avoided.  Finite-width energy--momentum kernels do not mean that closed systems violate conservation laws.  Rather, the finite kernel reflects that a finite effective transition description does not implement the asymptotic S-matrix idealization.

Exact conservation belongs to the full closed-system symmetry or to the asymptotic scattering limit.  Finite broadening belongs to finite observation, finite support, finite coherence, and finite admissible transition windows.

Thus the claim is:

\[
\text{finite window} \neq \text{fundamental nonconservation}.
\]

It is instead:

\[
\text{finite window} = \text{finite bookkeeping resolution}.
\]

\section{Scope of proof}

\begin{center}
\fbox{\begin{minipage}{0.88\linewidth}
\textbf{Proved here.}  Finite time-window kernels converge distributionally to
energy-conservation deltas.  The normalized squared sinc kernel yields the golden-rule rate
factor, and finite spacetime windows recover S-matrix deltas at infinite support.

\textbf{Conditional interpretation.}  These kernels may be read in MTT as finite bookkeeping
kernels only after a selected source derives the relevant transition window.

\textbf{Not proved here.}  Scattering theory, LSZ, Haag--Ruelle asymptotics, wave packets,
conservation laws, and an MTT window source remain external.
\end{minipage}}
\end{center}

\section{Conclusion}
\enlargethispage{2\baselineskip}

Finite-time scattering makes the sharp-limit structure explicit: a finite process produces a
finite kernel, while the Dirac delta appears in the asymptotic idealization.

The central result is:
\[
\boxed{
2\pi\delta(\Delta E)
=
\lim_{T\to\infty}
\int_{-T/2}^{T/2}e^{i\Delta E t}\,dt
}
\]
in the distributional amplitude sense, and
\[
\boxed{
\delta(\Delta E)
=
\lim_{T\to\infty}
\frac{1}{2\pi T}\left|
\int_{-T/2}^{T/2}e^{i\Delta E t}\,dt
\right|^2
}
\]
in the normalized transition-rate sense.

Thus S-matrix deltas are the asymptotic endpoints of standard finite transition kernels.  An
MTT preparation, detector, or fundamental-width source remains to be derived.

\begin{thebibliography}{9}

\bibitem{LSZ1955}
H. Lehmann, K. Symanzik, and W. Zimmermann,
``On the formulation of quantized field theories,''
\emph{Il Nuovo Cimento} \textbf{1} (1955), 205--225.
\url{https://doi.org/10.1007/BF02731765}

\bibitem{Ruelle1962}
D. Ruelle,
``On the asymptotic condition in quantum field theory,''
\emph{Helvetica Physica Acta} \textbf{35} (1962), 147--163.
\url{https://www.e-periodica.ch/digbib/view?pid=hpa-001:1962:35::472}

\end{thebibliography}

\end{document}
